\section{Vision Benchmarks, Metrics \& Evaluation}\label{sec:eval}
\lab{ImageNet-C \textperiodcentered\ COCO mAP \textperiodcentered\ Panoptic PQ \textperiodcentered\ FID 50k \textperiodcentered\ PSNR / SSIM / LPIPS}

\diagpair{%
\begin{tikzpicture}[x=1mm,y=1mm,every node/.style={font=\scriptsize}]
\node[box,draw=acc,fill=soft] (r) at (0,10) {Real Data\\$\mathcal{N}(\mu_r, \Sigma_r)$};
\node[box,draw=acc2,fill=soft2] (g) at (28,10) {Generated\\$\mathcal{N}(\mu_g, \Sigma_g)$};
\draw[arr,acc3,thick] (r) to[bend left=25] node[midway,above] {Inception-v3 pool3} (g);
\node[box,draw=acc3,fill=acc3!8] at (14,-1) {$\text{FID} = \|\mu_r - \mu_g\|^2 + \text{Tr}(\dots)$};
\draw[arr] (14,6) -- (14,2);
\end{tikzpicture}
}{%
\textbf{The Evaluation Landscape.} Visual evaluation spans categorical recognition, dense pixel geometry, perceptual fidelity, and multimodal reasoning.\\[2pt]
{\color{acc2}$\blacktriangleright$} \textbf{COCO $\text{mAP}$}: Evaluated across 10 IoU thresholds from $0.50$ to $0.95$ with $0.05$ step size.\\
{\color{acc2}$\blacktriangleright$} \textbf{FID}: Measures 2-Wasserstein distance between multivariate Gaussian features in Inception space.
}

\begin{mem}[Key Vision Benchmarks \& Thresholds]
\begin{tabularx}{\linewidth}{@{}lLL@{}}
\toprule
\textbf{Domain} & \textbf{Benchmark} & \textbf{Standard Metric}\\
\midrule
\textbf{Classification} & ImageNet-1K, ImageNet-A/R/C & Top-1 Acc ($\%$) / Mean Corruption Error (mCE)\\
\textbf{Detection} & MS COCO (80 classes) & $\text{mAP}_{50:95}$, $\text{AP}_{S}$, $\text{AP}_{M}$, $\text{AP}_{L}$\\
\textbf{Segmentation} & ADE20K (semantic), Cityscapes & Mean IoU (mIoU), Panoptic Quality ($\text{PQ} = \text{SQ} \times \text{RQ}$)\\
\textbf{VLM Reasoning} & MMBench, MathVista, DocVQA & Overall Accuracy, OCR-F1 score\\
\textbf{Hallucination} & POPE (Random, Popular, Adversarial) & Binary F1-score \& Yes-ratio\\
\textbf{Generative} & ImageNet $512^2$, COCO 30k & $\text{FID}$ ($<3.0$ frontier), CLIP Score ($>30.0$)\\
\textbf{3D Quality} & NeRF Synthetic, Tanks \& Temples & PSNR ($>32\text{ dB}$), SSIM ($>0.95$), LPIPS ($<0.05$)\\
\bottomrule
\end{tabularx}
\end{mem}

\begin{qa}[Metric Formulations \& Evaluator Pitfalls]
\Q{State the Fréchet Inception Distance (FID) formula and explain its vulnerabilities.}{
Given feature representations from the 2048-dimensional pool3 layer of Inception-v3:
$$\text{FID} = \|\mu_r - \mu_g\|_2^2 + \text{Tr}\left(\mathbf{\Sigma}_r + \mathbf{\Sigma}_g - 2(\mathbf{\Sigma}_r \mathbf{\Sigma}_g)^{1/2}\right)$$
This is the 2-Wasserstein distance between two continuous Gaussians $\mathcal{N}(\mu_r, \Sigma_r)$ and $\mathcal{N}(\mu_g, \Sigma_g)$.
\textbf{Vulnerabilities}:
\begin{enumerate}[leftmargin=*,itemsep=0.8pt]
\item \textbf{Sample Size Sensitivity}: FID is biased upward with small sample counts; comparing FID requires evaluating on exactly 50,000 generated samples ($N=50\text{k}$).
\item \textbf{Inception Bias}: Inception-v3 was trained on ImageNet and is blind to human anatomy, facial distortion, text spelling, and high-frequency noise.
\item \textbf{JPEG and Resizing artifacts}: Compressing images or using bilinear rather than bicubic resizing alters FID by up to $\pm 5$ points without perceptible image differences.
\end{enumerate}}
\Q{PSNR vs SSIM vs LPIPS: what does each measure, and when does each lie?}{
\textbf{PSNR} $= 10 \log_{10}\frac{\text{MAX}^2}{\text{MSE}}$ dB. Pixel-wise; a 1-pixel shift or slight blur can cost several dB while looking identical, and blurry averages score well. $+1$ dB $\approx$ 21\% lower MSE.\\
\textbf{SSIM} compares local means, variances and covariance (luminance, contrast, structure) in $11\times 11$ Gaussian windows:
$$\text{SSIM} = \frac{(2\mu_x\mu_y + C_1)(2\sigma_{xy} + C_2)}{(\mu_x^2 + \mu_y^2 + C_1)(\sigma_x^2 + \sigma_y^2 + C_2)}$$
Better on structure, still insensitive to texture realism.\\
\textbf{LPIPS}: distance between unit-normalized deep features (AlexNet/VGG) with learned per-channel weights, trained on human 2AFC judgements. Tracks perceived similarity best; lower is better.\\
\textbf{Rule}: report all three for reconstruction (NeRF/3DGS, super-resolution, VAE); a GAN-sharpened output can lose PSNR but win LPIPS. For generation without a reference image, use FID/CLIPScore instead.}
\Q{How is Panoptic Quality (PQ) decomposed into SQ and RQ?}{
Kirillov et al. decompose Panoptic Quality as $\text{PQ} = \text{SQ} \times \text{RQ}$:
$$\text{SQ} = \frac{\sum_{(p, g) \in TP} \text{IoU}(p, g)}{|TP|}, \quad \text{RQ} = \frac{|TP|}{|TP| + \frac{1}{2}|FP| + \frac{1}{2}|FN|}$$
\begin{itemize}[leftmargin=*,itemsep=0.8pt]
\item \textbf{Segmentation Quality (SQ)}: Measures average boundary overlap purely among true-positive matched segments ($\text{IoU} > 0.5$).
\item \textbf{Recognition Quality (RQ)}: Standard F1-score measuring detection/classification accuracy of segments regardless of pixel tightness.
\end{itemize}
Decouples mask boundary quality from false-positive detection errors.}
\end{qa}

\begin{trap}[Benchmarking Traps]
\begin{itemize}
\item \textbf{ImageNet-C Corruption Leakage}: Using Gaussian blur or noise during training data augmentation contaminates ImageNet-C evaluation, invalidating reported robustness metrics.
\item \textbf{CLIPScore Gamification}: Generating images with giant repeated text labels exploits CLIP's text-matching bias, yielding high CLIPScore despite awful visual image quality.
\end{itemize}
\end{trap}
